\section{Results}
\label{sec:results}

\tbd{Every number in this section comes from a draft run (corridor scale 8, six departures, one round of rig sizing) on code that predates the latest changes to the router, the kite's inputs and the speed steps. The final run replaces them. That draft ran e-methanol where this paper now runs e-ammonia, so its numbers on the e-fuel stand in for those of the next run.} Throughout, a cost saving is the drop in cost per tonne-kilometre that a rig gives against the same ship without a rig on the same fuel and routes, and a point is one percentage point of that saving. Every saving assumes that each rig reaches its goals (\cref{sec:res_goals}), and \cref{sec:certainty} asks how sure the comparison is. Figures on the scaled family of ships, one \SI{85}{m} coaster hull scaled geometrically to lengths of up to \SI{200}{m} (\cref{sec:ships}), on the route to Guadeloupe show how a result comes about. The figures over the thirteen classes and six routes carry the results themselves.

\subsection{Which rig where, and how much cost it saves}
\label{sec:res_fuel}
\label{sec:res_routes}

\Cref{fig:class_comparison} gives the cost saving of the best size of each rig on every class, from an \SI{85}{m} coaster to a \num{24000}~TEU Megamax container ship (\cref{tab:classes}). On VLSFO, the fossil fuel oil of today's ships, at today's prices the best rig saves 1.8\,\% of the cost on the coaster to 5.5\,\% on the Kamsarmax on average over the routes. The cost saving does not grow with the size of the ship. Bulk carriers and tankers save 4.7 to 5.5\,\% and container ships 2.7 to 3.7\,\%. Wing sails save at most 1.4\,\% of the cost on any class, and they cost more than they save on the coaster and the three smaller container ships. On e-ammonia the kite saves 12.2\,\% of the cost on average over all classes and routes, the rotors 11.4\,\% and the wing sails 6.3\,\%. The gap between kite and rotors is within the error of this draft. The rotors lead on the car carrier, 12.5 against 10.7\,\%, and the kite on the four container ships, which sail fastest.

\Cref{fig:winner_map} names the rig that lowers the cost most on each of the 78 combinations of class and route. On e-ammonia the kite leads in 48 and the rotors in 30. On VLSFO the kite leads in 66 and the rotors in 11, and on the coaster to Bilbao no rig pays. Wing sails lead in none. On the container ships and the car carrier the class decides. On e-ammonia on the bulk carriers and tankers kite and rotors are within a point of each other, at 12 to 14\,\%, and the route decides. The kite leads to Guadeloupe, Houston and Bilbao on nearly every class and the rotors on most classes to Dakar and Tubarão. The route also spreads the cost saving more than the class does. On e-ammonia the kite saves 15.2\,\% of the cost on average to Guadeloupe and the rotors 16.2\,\% to Tubarão. The better of the two saves 11.6\,\% of the cost to Dakar and only 3.8 to 5.9\,\% on the short route to Bilbao, on which the kite flies 25 to 29\,\% of the time instead of 41 to 61\,\%. \Cref{sec:res_forces} explains this pattern from where each rig pulls.

\begin{figure*}[tbp]
\centering
\plot[\textwidth]{class_comparison}
\caption{Cost saving of the best size of each rig for the thirteen ship classes with a VLSFO (a) and an e-ammonia (b) drive: the mean over the six routes (marker) and the range over them (bar).}
\label{fig:class_comparison}
\end{figure*}

\begin{figure*}[tbp]
\centering
\plot[\textwidth]{winner_map}
\caption{The rig that lowers the cost per tonne-kilometre most for each ship class and route, with a VLSFO (a) and an e-ammonia (b) drive. The number is its cost saving in percent and the shade grows with it. In grey cells no rig pays and the number is the smallest loss.}
\label{fig:winner_map}
\end{figure*}

The rigs are sized for cost, so the fuel they save is larger than their cost saving (\cref{tab:fuel}). On e-ammonia the cheapest size of the rotors saves 29 to 39\,\% of the fuel of each class, of the kite 30 to 38\,\% and of the wing sails 14 to 29\,\%. We compare on cost and fuel energy and give CO$_2$ only for VLSFO, where the CO$_2$ saved is the fuel saved times a fixed factor. E-ammonia holds no carbon, so on that drive only the VLSFO pilot emits CO$_2$ tank-to-wake, that is, in what leaves the ship's stack, 5\,\% of what the same energy of VLSFO would. A well-to-wake account, which also counts the emissions of making the fuel, depends on how its hydrogen is produced and on its nitrous oxide and slip on the engine. These are outside this comparison.

\begin{table*}[tbp]
\centering
\footnotesize
\caption{Fuel energy saved by the cheapest size of each rig on each ship class with each drive, the mean over the six routes, in percent of the fuel of the same ship without a rig, and the tank-to-wake CO$_2$ of each class on VLSFO without a rig. On VLSFO the CO$_2$ saved equals the fuel saved.}
\label{tab:fuel}
\InputIfFileExists{sections/fuel_table}{}{\tbd{Fuel table pending: run the fuel\_table step.}}
\end{table*}

Part of each cost saving is time, since a rig makes speed cheaper and the router sets speed by cost (\cref{sec:res_speed}). \Cref{tab:split} splits the cost saving of the cheapest size of each rig into three parts. The first is the fuel it saves. The second is the time it saves: the crew, capital and cargo interest (the interest on the value of the cargo while it is on board) of the ship without a rig over the hours the rig takes off the round trip. The third is the rig's own capital, maintenance and wear and the cargo its mass displaces. On e-ammonia the kite and the rotors gain 3.4 to 4.8 points of cost from time, against 9.9 to 12.8 from fuel, and the wing sails 1.7 to 1.9 points from time. On VLSFO, with smaller rigs, time adds at most 0.5 points. The fuel saved alone pays for the rig everywhere except for the wing sails on VLSFO, whose 0.2\,\% cost saving comes from time.

\begin{table*}[tbp]
\centering
\footnotesize
\caption{Cost saving of the cheapest size of each rig with each drive, the mean over the classes and routes, split into the fuel it saves, the time it saves at the crew, capital and cargo interest of the ship without a rig, and the rig's own cost and the cargo it displaces, each in points of the cost of the ship without a rig; and the share of round trips the rig shortens and the median change of their duration.}
\label{tab:split}
\InputIfFileExists{sections/split_table}{}{\tbd{Split table pending: run the split\_table step.}}
\end{table*}

\Cref{fig:premium} shows how much of the cost premium of an e-ammonia ship, its extra cost per tonne-kilometre over the same ship on VLSFO without a rig, each rig closes. Without a rig the premium is 56 to 60\,\% on the scaled ships. The kite closes 31\,\% of it on the \SI{85}{m} ship and 40\,\% on the larger ones, the rotors 27\,\% on the \SI{85}{m} ship to 42\,\% on the \SI{200}{m} ship and the wing sails 14 to 19\,\%. With the better of kite and rotors the e-ammonia ship still costs 32 to 38\,\% more per tonne-kilometre than the same ship on VLSFO without a rig.

\begin{figure}[tbp]
\centering
\plot{premium}
\caption{Share of the cost premium of an e-ammonia ship over the same ship on VLSFO without a rig that the best size of each rig closes, over waterline length on the route to Guadeloupe. At 100\,\% a ship on that fuel with that rig costs as much per tonne-kilometre as a fossil-fuelled ship today.}
\label{fig:premium}
\end{figure}

Sizing a rig for the lowest cost leaves fuel unsaved that a larger rig would save at a small extra cost. \Cref{fig:abatement} traces cost against fuel saved as each rig grows to the largest that fits. Where a curve returns to the cost of the same ship without a rig, the rig saves the most fuel it can at no extra cost: \tbd{fuel saved at break-even per rig and drive}. On VLSFO the slope beyond that point is the cost of the CO$_2$ abated, \tbd{USD per tonne}. \Cref{tab:payback} gives the discounted payback time of the cheapest rig: the years until its yearly saving in running cost, before its own capital charges, repays its price at the discount rate (\cref{sec:cost_model}). On e-ammonia the kite pays back in 0.5 to 1.5 years, the rotors in 1.4 to 4.1 and the wing sails in 1.9 to 5.5. On VLSFO the kite takes 1.5 to 5.1 years, the rotors 4.1 to 15 and the wing sails 6.8 to 23. On three classes the wing sails never pay back.

\begin{figure*}[tbp]
\centering
\plot[\textwidth]{abatement}
\caption{Cost per tonne-kilometre against the fuel energy saved, relative to the same ship on VLSFO without a rig, for every size of each rig on the \SI{140}{m} ship to Guadeloupe (top) and a Kamsarmax bulk carrier to Tubarão (bottom), one column per drive. The black point is the ship without a rig on that drive. Each line joins the sizes of one rig in order of area, from the smallest (open marker) to the largest that fits on deck. Where a larger kite flies a smaller share of the time, it saves less fuel and the line turns back.}
\label{fig:abatement}
\end{figure*}

\begin{table*}[tbp]
\centering
\footnotesize
\caption{Discounted payback time in years of the cheapest size of each rig on each ship class, the median over the six routes, at the base discount rate. A dash marks where the median rig does not pay back within its life.}
\label{tab:payback}
\InputIfFileExists{sections/payback_table}{}{\tbd{Payback table pending: run the payback\_table step.}}
\end{table*}

\subsection{Another e-fuel}
\label{sec:res_fuels}

The comparison rests on one e-fuel (\cref{sec:fuel_choice}). What sets a rig's value on a ship is mostly the price of the energy it saves, so \cref{fig:fuels} sweeps that price for e-ammonia from \SIrange{30}{120}{USD/GJ}. It also runs the same ships on e-methanol from \SIrange{40}{120}{USD/GJ}, as that fuel would be installed, with a diesel-electric drive and a tank that takes no cargo space. One assessment puts 2030 prices inside that sweep: \SIrange{47}{82}{USD/GJ} for e-ammonia, \SIrange{73}{118}{USD/GJ} for e-methanol and \SIrange{60}{98}{USD/GJ} for e-LNG \citep{lrumas2020}. \tbd{At the same price the two e-fuels give cost savings within X points of each other on both ships, and the rig that saves most is the same on both at every price swept; the lead changes only below Y USD/GJ, where the rotors pass the kite on the Kamsarmax.} An e-fuel differs from another mainly in what the engine and the tank cost, which every rig shares, and in the tank space and mass a rig saves, largest for ammonia. E-LNG is not in the model. Its price lies between the other two, so the sweep covers it, but its insulated tank would give a rig that saves fuel some more to gain, as ammonia's tank does.

\begin{figure*}[tbp]
\centering
\plot[\textwidth]{fuels}
\caption{Cost saving of the best size of each rig over the price of the e-fuel on the \SI{140}{m} ship to Guadeloupe (a) and a Kamsarmax bulk carrier to Tubarão (b), on e-ammonia in a dual-fuel engine (solid) and on e-methanol in a diesel-electric drive (dashed), from the sizes of \cref{fig:situations}.}
\label{fig:fuels}
\end{figure*}

\subsection{Why the rigs differ}
\label{sec:res_forces}

Before any voyage, the three rigs differ in the force they put on the ship at each wind direction. \Cref{fig:rig_polar} shows the thrust, the force along the ship's course, and the side force, the force across it, on a \SI{140}{m} ship at \SI{10.5}{kn} in a \SI{10}{m/s} wind, each rig at the setting that needs the least engine power. The direction of the wind is the true wind angle, measured from the bow. In sailing terms the ship is close-hauled with the wind from ahead at a small angle, on a beam reach with the wind from the side, on a broad reach with it from behind the beam, and dead downwind with it from astern. A rotor at its working spin ratio, the speed of its surface over the wind speed, has a lift coefficient of 9.5 and a wing at most 2.0 (\cref{sec:deck_rigs}), so a square metre of rotor pulls about five times as hard as a square metre of wing. The rotors pull hardest from close-hauled to beam reach and the kite from a broad reach to dead downwind. The wing sails pull less than the rotors at every angle\tbd{ and more than the kite below about \SI{55}{\degree}}. Side force follows the lift. The hull resists it by moving slightly sideways through the water, its leeway, and the rudder holds the course against any turning moment the rig adds. Both add drag.

\begin{figure}[tbp]
\centering
\plot{rig_polar}
\caption{Thrust (top) and side force (bottom) over the true wind angle of an \SI{800}{m^2} kite, \SI{800}{m^2} of wing sails and \SI{400}{m^2} of rotors on a \SI{140}{m} ship at \SI{10.5}{kn} in a \SI{10}{m/s} wind at \SI{10}{m}, each at the setting that needs the least engine power.}
\label{fig:rig_polar}
\end{figure}

On a voyage the router turns this thrust into fuel or speed. \Cref{fig:passage} follows one winter passage to Guadeloupe. The ships without a rig and with wing sails stay close to the great circle, the shortest path over the globe. The kite and rotor ships route further south into the trade winds, the steady easterly winds of the subtropics. The kite pulls in long bursts. While it flies, its ship speeds up to as much as 16~kn and its fuel power, the power the engines draw from fuel, drops to almost nothing. Between flights the ship sails on the engine. The rotors pull more steadily and less, and their ship sails at 10 to 13~kn for much of the passage. Both arrive about two days before the ship without a rig. The wing-sail ship keeps to the speed of the ship without a rig.

\begin{figure*}[tbp]
\centering
\plot[\textwidth]{passage}
\caption{One outbound passage to Guadeloupe, departing on 1 January 2017, of a \SI{140}{m} e-ammonia ship without a rig and with a \SI{3200}{m^2} kite, \SI{1500}{m^2} of wing sails or \SI{800}{m^2} of rotors: (a) the routed tracks with the great circle dotted, over the \SI{10}{m} wind six days after departure, each track narrower than the one before it in the legend so that tracks which coincide stay visible, (b) speed through the water, (c) rig thrust and (d) fuel power over the passage, as 12-hour moving averages.}
\label{fig:passage}
\end{figure*}

\Cref{fig:wind_angle} shows where on a voyage each rig earns its cost saving, from the routed round trips of the \SI{140}{m} e-ammonia ship with a \SI{3200}{m^2} kite, \SI{1500}{m^2} of wing sails and \SI{800}{m^2} of rotors. Without a rig the ship spends \tbd{28}\,\% of its time at sea close-hauled, below \SI{60}{\degree}, \tbd{28}\,\% abeam and \tbd{44}\,\% downwind, beyond \SI{120}{\degree}, but the routes differ. The kite earns \tbd{70}\,\% of its propulsive work, thrust times distance, downwind, and the rotors and wing sails \tbd{42 to 47}\,\% abeam and about as much downwind. Per dollar of rig cost, kite wear included, the kite gives the most thrust on every route, \tbd{50 to 116}\,kN per USD~\num{1000} a day, against \tbd{21 to 62} for the rotors and \tbd{18 to 45} for the wing sails. \tbd{Explain the route pattern of \cref{sec:res_fuel} from these wind angles once recomputed.} \Cref{fig:route_comparison} shows the cost saving per route on three ships of the scaled family at two fixed sizes of each rig.

\begin{figure*}[tbp]
\centering
\plot[\textwidth]{wind_angle}
\caption{Where the wind comes from and where each rig pulls, from the weather-routed round trips of the \SI{140}{m} e-ammonia ship over the sampled departures on every route: (a) share of the time at sea without a rig in each band of true wind angle, outbound (left bar) and homeward (right bar); (b) mean thrust of a \SI{3200}{m^2} kite, \SI{1500}{m^2} of wing sails and \SI{800}{m^2} of rotors in each band over all routes, each on its own routed track; (c) share of each rig's propulsive work, thrust times distance, from each band per route; (d) mean thrust at sea over the rig's cost per day at sea: the fixed cost it adds to the ship, spread over the round trip with its port calls, plus canopy and tether wear and launches of the kite.}
\label{fig:wind_angle}
\end{figure*}

\begin{figure*}[tbp]
\centering
\plot[\textwidth]{route_comparison}
\caption{Cost saving of the better of two sizes of each rig against the same ship without a rig, per route, for ships of \SIlist{85;140;200}{m} with a VLSFO (top) and an e-ammonia (bottom) drive: kites of \SIlist{800;3200}{m^2}, wing sails of \SIlist{800;1500}{m^2} and rotors of \SIlist{400;800}{m^2}, of which only \SI{400}{m^2} fits on the \SI{85}{m} ship. On the \SI{200}{m} ship the rotors are smaller than their best size of \cref{fig:rig_comparison}.}
\label{fig:route_comparison}
\end{figure*}

Both deck rigs stand on the same deck and are priced on the same basis with the same life of 25 years (\cref{sec:costs,sec:inputs}), so the gap between them shows what the rotor's higher lift is worth. A square metre of rotor costs about three times a square metre of wing and pulls about five times as hard. In \cref{fig:breakdown} wing sails and rotors of similar thrust add about the same capital and maintenance per tonne-kilometre. The difference is the fuel. On e-ammonia the rotors save nearly twice what the wings save (\cref{fig:class_comparison}). On VLSFO fuel is a quarter of the cost of the \SI{85}{m} ship, too small a part for either deck rig to pay. The smallest wing rig saves USD~0.21 per \num{1000} tonne-kilometres and adds USD~0.42 of capital and maintenance.

\begin{figure}[tbp]
\centering
\plot{breakdown}
\caption{Cost per tonne-kilometre, in USD per \num{1000}~tkm, split into its parts on the route to Guadeloupe for an \SI{85}{m} and a \SI{200}{m} ship on VLSFO without a rig, and on e-ammonia without a rig and with the cheapest size of each rig, whose area is given with each bar. Capital includes the fuel store.}
\label{fig:breakdown}
\end{figure}

\subsection{Rig size}
\label{sec:res_size}

\Cref{fig:class_sizes} gives the rig area that lowers the cost most on each class, and \cref{fig:rig_comparison} how the cost saving changes with area on the scaled family. On most large ships the best wing sails and rotors are the largest that fit on deck, with up to 15 wings and 23 rotors on the Megamax container ship. On these ships the deck, a physical limit, sets their size and their cost saving. The best kite reaches \SI{16000}{m^2}, the upper bound of the search, on the Megamax container ship on every route on e-ammonia and on five of the six on VLSFO. It reaches it on the Neo-Panamax container ship on four or five of the six routes and on the VLCC on two to four. The largest ship kite designed so far is \SI{1000}{m^2} \citep{airseas_seawing_web}. The bound is a limit of the model, so on these three ships the kite's best size and cost saving are not known, only that they are at least this large. At the other end the best kite of the coaster is the smallest searched, \SI{500}{m^2}, on two or three of the six routes. On the \SI{140}{m} e-ammonia ship to Guadeloupe the kite saves \tbd{14.9}\,\% of the cost at \SI{1600}{m^2} and about two points less for each doubling or halving of its area. Half the best rotor area costs that ship 2.4 points and half the best wing area 1.9 points. \Cref{sec:disc_size} takes up what these sizes mean.

\begin{figure*}[tbp]
\centering
\plot[\textwidth]{class_sizes}
\caption{The rig area that lowers the cost per tonne-kilometre most for each ship class with a VLSFO (a) and an e-ammonia (b) drive: the geometric mean over the six routes (marker) and the range over them (bar). The number next to the wings and rotors is the median number of units at that size. Kites were searched from \SIrange{500}{16000}{m^2}, wings and rotors from a fifth to all of what fits on deck. Where a bound was the cheapest size searched, the search also tried the size a factor $2^{1/4}$ inside it.}
\label{fig:class_sizes}
\end{figure*}

\begin{figure*}[tbp]
\centering
\plot[\textwidth]{rig_comparison}
\caption{Cost saving against the same ship without a rig on the route to Guadeloupe, for every size of kite, wing sails and rotors (rows) and ship length (columns), with a VLSFO (a) and an e-ammonia (b) drive. The outlined cell is the best size of each rig on each ship. Blank cells are sizes that do not fit on deck or were not run. Wings and rotors were run on three of the five lengths.}
\label{fig:rig_comparison}
\end{figure*}

\subsection{How a rig changes the ship's speed}
\label{sec:res_speed}

Ships with a rig sail faster than the same ship without one, because the router sets speed by cost and not by a schedule. Over the 36 cases of ship length, drive and route, each with its cheapest rig size, the round trip of the kite ship is shorter by a median of \tbd{6.3}\,\%, faster in \tbd{X} cases and slower in \tbd{Y}. Wing sails shorten it by \tbd{3.7}\,\% and rotors by \tbd{10.8}\,\%, and both are faster in all 36 cases. The effect follows the fuel price. The kite gains \tbd{1.1}\,\% on VLSFO and \tbd{10.0}\,\% on e-ammonia, the rotors \tbd{5.5} and \tbd{13.1}\,\%. Dear fuel makes the ship without a rig steam slowly, so a rig has more speed to give back.

A rig that gave the same thrust at every speed would not change the best speed. The cost per mile is $(R(V)-F_x)/\eta + c/V$ for a resistance $R$, an efficiency $\eta$ and a cost $c$ per hour of time, and a constant $F_x$ drops out of its derivative. \Cref{tab:speed_choice} shows what moves the best speed instead. The kite moves it by the speed its pull alone can reach. Up to there the ship burns only its hotel load, so sailing faster costs nothing and saves time. Holding the kite's force at its value at the slower speed barely changes the result. The rotors' best speed rises mainly because their force grows with speed, and with the force held it stays close to that of the ship without a rig. The wing sails move it by a knot or less, mostly by lightening the load on the propeller. The capital and wear of the rig, which also make time dearer, shift the best speed by about \tbd{0.1}~kn.

\begin{table}[tbp]
\centering
\caption{Cost-optimal speed [kn] of the \SI{140}{m} ship in steady wind, median over winds of 6, 10 and \SI{14}{m/s} at 60, 90, 120 and \SI{150}{\degree}, with the reference rigs: with the rig's capital and wear left out of the cost, with its force held at its value at the speed of the ship without a rig, and the speed at which that held force alone moves the ship. \tbd{Draft run.}}
\label{tab:speed_choice}
\adjustbox{max width=\columnwidth}{\InputIfFileExists{sections/speed_choice_table}{}{\tbd{Speed table pending: run the speed\_choice\_table step.}}}
\end{table}

\Cref{tab:speed_terms} puts the same ships on the terms other studies use. At a fixed speed the kite on VLSFO saves \tbd{5.9}\,\% of the fuel at \SI{8}{kn}, \tbd{9.5}\,\% at \SI{9}{kn} and \tbd{19.0}\,\% at \SI{10}{kn}, because at a low speed the fuel it can save is worth little more than its wear and the router keeps it down. On e-ammonia it saves \tbd{21.9}, \tbd{26.2} and \tbd{23.0}\,\% of the fuel. In the time of the ship without a rig it saves \tbd{35.9}\,\% of the fuel on VLSFO and \tbd{44.0}\,\% on e-ammonia, against the \tbd{31.0} and \tbd{37.0}\,\% it saves when cost sets the speed. The speed is bought with part of the fuel saving.

\begin{table*}[tbp]
\centering
\caption{Fuel saved and time saved by the reference rigs on the \SI{140}{m} ship on two routes, against the ship without a rig on the same terms: at a fixed speed, in the same time as the ship without a rig routed for cost, and with speed set by cost. A fixed-speed ship drops to the highest of its routing speeds it can hold where it cannot hold that speed, which a rig does more often, so its time changes too. \tbd{Draft run.}}
\label{tab:speed_terms}
\InputIfFileExists{sections/speed_terms_table}{}{\tbd{Terms table pending: run the speed\_terms\_table step.}}
\end{table*}

\Cref{tab:flight} shows where the kite flies. Upwind of \SI{30}{\degree} it could pull for only \tbd{13}\,\% of the time and is rarely up. From \SIrange{60}{150}{\degree} it flies for most of the time it could pull, and the router keeps it down where the thrust would not pay for its wear. It stays up through some lulls, because landing and launching cost more than waiting, and \tbd{17}\,\% of its hours aloft give no thrust. Over all 36 cases the kite is up for \tbd{45}\,\% of the time. The lowest apparent wind angle at which it pulls is \tbd{14.9}\,\si{\degree}, against the limit of \cref{eq:upwind_limit}. Wing sails and rotors pull down to \SI{6.4}{\degree} and \SI{8.5}{\degree}, because their lift to drag is higher than the kite's whole-system value. Tests hold all three rigs to their limits. The kite is up for more of the hours at higher ship speeds, because the router raises the speed where the kite pulls hardest. On VLSFO the kite ship also spends more of its time below \SI{10}{kn}, \tbd{25} against \tbd{16}\,\%, since it can idle in the headwinds and make up the time when the wind helps.

\begin{table*}[tbp]
\centering
\caption{Where the kite works over the 36 cases, each with its cheapest size, by band of true wind angle: share of the time in the band, share of that time in which the kite could pull, is up and pulls, its mean thrust while up, lulls and upwind drag included, the share of the hours at each ship speed in which it could pull and in which it is up, the speed of the ship without a rig and of the kite ship in that band, and the lowest apparent wind angle (AWA) at which each rig pulls. \tbd{Draft run.}}
\label{tab:flight}
\InputIfFileExists{sections/flight_table}{}{\tbd{Flight table pending: run the flight\_table step.}}
\end{table*}

Published kite studies report fuel savings at given speeds and with the kite's size fixed, so we compare on those terms (\cref{app:speed_literature} gives their statements in full). \citet{cedelft2016}, at 100 and 85\,\% of service speed, find that kites save relatively more at the lower speed and conclude that they are best suited to lower ship speeds. \citet{emsa2023} state the opposite, on one provider-reported pair of figures. At a higher ship speed the apparent wind moves ahead and the kite's force falls. \citet{podeur2016} find that going from 16 to \SI{19}{kn} cuts the round-trip fuel saving from 12.1 to 6.6\,\%, with less flight time. By our calculation from their tables, most of that drop is the larger fuel use of the ship at the higher speed, which rises by 42\,\%, while the absolute fuel saving falls by 22\,\%.

\Cref{tab:published} sets these results beside our model at the same speeds on the nearest of our ships. In every study, as in ours, a kite of one size saves a smaller share of the fuel at a higher speed. Our kite saves less fuel than every single published figure, from half or less of the routed range of \citet{bentin2018} to under a tenth of the single crossing of \citet{leloup2013}. Several differences each raise the published figures, and we have not measured the share of each (\cref{app:speed_literature}): a higher glide ratio \citep{podeur2016,leloup2013}, no cost for the side force, and single crossings against our average over both directions and six departures. The ranges of CE Delft and \citet{traut2014} bracket ours. They span ships and routes and not the conditions of one ship. Studies that do let the speed vary work with rotors and minimise fuel under a deadline \citep{mason2023,dhainaut2025}, or optimise speed without a rig \citep{tillig2020b}. Here the speed varies with a rig under an objective that values time. A kite that delivers less at a higher speed is compatible with a best speed that rises, since the best speed is set by where the rig's pull runs out and by the value of time.

A slower ship gives a larger share to a kite of one size, but so does a larger kite, so a share at one size does not show which speed suits kites. \Cref{fig:speed_size} tests this at held speeds on the \SI{140}{m} ship and the feeder container ship, with each kite flown once free of its running cost, as in the published studies, and once priced as in the rest of this paper (\cref{sec:further_speed}). At one size the share of fuel saved falls with speed, as published, and most of that fall is the larger fuel use it is divided by. The \SI{400}{m^2} kite on the e-ammonia ship saves \tbd{11.7}\,\% of the fuel at \SI{8}{kn} and \tbd{9.8}\,\% at \SI{11}{kn}, while the ship's fuel grows by \tbd{38 to 43}\,\% and the energy the kite saves by \tbd{20}\,\%.

A larger kite takes the share back. On the \SI{140}{m} ship the kite that saves the most fuel saves a larger share at \SI{11}{kn} than at \SI{8}{kn}, \tbd{27.7} against \tbd{23.2}\,\% on e-ammonia. Priced on e-ammonia, every size saves more of the cost at the higher speed. On VLSFO the cheapest kite saves at most \tbd{0.9}\,\% of the cost at any of these speeds.

\begin{figure*}[tbp]
\centering
\plot[\textwidth]{speed_size}
\caption{Fuel and cost savings of kites at a held ship speed for each kite size, on both routes, against the same ship without a rig at the same speed: fuel saved by the kite free of its running cost, in percent (a, d, g) and in GJ per round trip (b, e, h), and cost per tonne-kilometre saved by the kite priced as in the rest of this paper (c, f, i), on the \SI{140}{m} ship on VLSFO (a--c) and e-ammonia (d--f) and on the feeder container ship on VLSFO (g--i). \tbd{Draft run.}}
\label{fig:speed_size}
\end{figure*}

On the feeder from 13 to \SI{19}{kn} this holds only in part. The \SI{12800}{m^2} kite does not buy the share back. Its weight, its tether's force limit and the larger wind it needs to fly hold it down, and we have not separated these. The kite's price is the second limit. The kite that saves the most cost is two to eight times smaller than the one that saves the most fuel, so the size that would take back the share is bought only where fuel is dear.

\begin{table}[tbp]
\centering
\caption{Published fuel savings of ship kites at fixed speed beside our model at the same speed on the nearest of our ships and on the route to New York: with the kite priced as in the rest of this paper, and free of running cost as in the published studies. Ships, routes, weather and kite models differ, so a row compares terms and not one experiment. \tbd{Draft run.}}
\label{tab:published}
\adjustbox{max width=\columnwidth}{\InputIfFileExists{sections/published_table}{}{\tbd{Published table pending: run the published\_table step.}}}
\end{table}

The speed result has limits. The router knows the weather ahead and no schedule binds it. A ship with fixed arrival times or port slots keeps the time of the ship without a rig and takes its gain as fuel (\cref{tab:speed_terms}). The gain depends on the value of time, which holds capital, crew and cargo interest (\cref{sec:cost_model}), and on the fuel price. Above the speed its engine alone could reach, a ship is sailing under its rig, and the \SI{140}{m} VLSFO ship reaches \SI{11}{kn} on its \SI{2.2}{MW}. We charge no drag for a propeller trailing at zero thrust, which flatters the fastest sailing hours, \tbd{3.7}\,\% of the kite's time above \SI{14}{kn}. The side force of the kite is balanced at each speed, but the hull's resistance at these speeds rests on the same calibration as at service speed (\cref{sec:disc_common}).

\subsection{VTOL kites are too expensive at ship scale}
\label{sec:res_rigid}

A VTOL kite is a rigid wing that carries propellers, which lift it off the deck and land it hovering like a helicopter, and flies loops on its tether in between (\cref{sec:rigid_kite}). It flies at ship scale but costs far more per unit of thrust than a soft kite (\cref{fig:rigid_scaling}). At \SI{12}{m/s} of apparent wind on a beam course the \SI{54}{m^2} wing costs USD~\num{16200} per kilonewton and a \SI{1000}{m^2} soft kite USD~\num{2360}, a seventh of that. The rigid wing's life of 15 years closes part of that gap. Over a year of \num{3000} flight hours and 150 launches, each kite at its goals, the wing costs USD~\num{2990} per kilonewton and the soft kite USD~790. \tbd{The prices per kilonewton over size that follow come from before the VTOL kite's drive lost the ground electronics that the ship's supply already holds, which lowers them by a few per cent, more for large wings; the final run replaces them.} Size helps at first and then hurts, because the wing's mass grows with its span to the power $\kappa$, the mass exponent (\cref{sec:rigid_kite}). At $\kappa=2.75$ the price per kilonewton falls to USD~\num{16500} at \SI{54}{m^2} and rises to USD~\num{34700} at \SI{500}{m^2}. The heavier wing also needs more wind to reach the top of the circle it flies on its tether (\cref{eq:loop_top}), \SI{6}{m/s} at \SI{30}{m^2} and \SI{12}{m/s} at \SI{500}{m^2}. The speed of the ship does not close this gap, since wind from ahead gives drag and not thrust.

\begin{figure*}[tbp]
\centering
\plot[\textwidth]{rigid_scaling}
\caption{VTOL kites and soft kites over size: (a) thrust of one kite on a beam course over the apparent wind, (b) price per thrust and (c) cost over a year of \num{3000} flight hours and 150 launches per thrust at \SI{12}{m/s} of apparent wind, for VTOL kites at three mass exponents and for soft kites, each at its goals. Wings that do not fly at \SI{12}{m/s} are left out.}
\label{fig:rigid_scaling}
\end{figure*}

If it reaches its goals, a VTOL kite does not pay on VLSFO at any size and only breaks even on e-ammonia (\cref{fig:rigid_kite}). On the \SI{140}{m} ship and the Kamsarmax bulk carrier its best size costs \tbd{2.7} and \tbd{2.0}\,\% more than no rig on VLSFO, on every route, and on e-ammonia it saves \tbd{$-0.3$} and \tbd{0.1}\,\% of the cost. That best size is nearly always one kite of \SI{30}{m^2}, the smallest searched, because it loses least (\cref{fig:rigid_size}), so the cheapest choice at these inputs is no rig.\tbd{ The routed numbers of this section are at $\kappa=2.5$, from before the move to 2.75. Between the runs at 2.5 and 3, that move adds 0.1 to 0.2 points to the best size and takes more from large wings; the final run replaces them.} It flies for \tbd{9}\,\% of the time on VLSFO and \tbd{23}\,\% on e-ammonia, against \tbd{26} and \tbd{44}\,\% for the soft kites of the same study on average. At USD~256 per launch and landing and USD~18 of upkeep per flight hour a small kite has to fly long or in strong wind to earn its flight. We read this as the reason but did not separate it from the data.

\begin{figure*}[tbp]
\centering
\plot[\textwidth]{rigid_kite}
\caption{Cost saving against the ship without a rig of the best size of the soft kite, wing sails, rotors and VTOL kites, at base inputs and with every input of the VTOL kite at its favourable bound, on the \SI{140}{m} ship and the Kamsarmax bulk carrier with a VLSFO and an e-ammonia drive. Bars are the mean over the six routes, whiskers the range over them. VTOL kites run from one to eight kites of \SIlist{30;54;100}{m^2} and one of \SIlist{200;300;500}{m^2}, as far as they fit on deck.}
\label{fig:rigid_kite}
\end{figure*}

\begin{figure*}[tbp]
\centering
\plot[\textwidth]{rigid_size}
\caption{Cost saving of VTOL kites against the ship without a rig over their total area, on the \SI{140}{m} ship (left) and the Kamsarmax bulk carrier (right) on VLSFO, as the mean over the six routes: one line for a single kite of any area and one each for fleets of kites of \SIlist{30;54;100}{m^2}, solid for the airframe mass of \cref{eq:rigid_mass} and dashed for a mass growing as the cube of the span.}
\label{fig:rigid_size}
\end{figure*}

Whether one wing or several comes out cheaper is decided only at sizes that do not pay. Up to \SI{100}{m^2} in total they tie, and at \SI{200}{m^2} several wings win by 2 to 10 points. The loss from several kites sharing one loop (\cref{eq:loop_share}) and the deck, station and tether that each kite brings hold several small wings back about as much as mass holds one large wing back.

With every input of the VTOL kite at the bound that favours it at once, it saves \tbd{2.3} and \tbd{2.1}\,\% of the cost on VLSFO and \tbd{9.0} and \tbd{9.4}\,\% on e-ammonia on the two ships, behind the soft kite and the rotors at base inputs (\cref{fig:rigid_kite}). Its best size is then mostly one or two kites of \SI{54}{m^2} or one of \SI{100}{m^2}. At today's canopy price it leads in \tbd{1} of the 24 combinations of ship, drive and route, \tbd{to recheck at the goal price}. At its goals the VTOL kite ranks last of the four rigs on both ships and both fuels. Adding it changes no leader among the rigs at base inputs, so it is left out of every other study. Its goals are a life of 15 years and the loss of one kite in \num{10000} launches and landings. A crash costs a whole airframe, USD~1.74 million for a \SI{54}{m^2} kite with half its tether, against USD~\num{233000} for a \SI{1600}{m^2} soft kite.
